% !TEX root = ../main.tex
\begin{ouverture}[gravitation]{11}{Interactions fondamentales et champs}{Partie 2 · Mouvement et interactions}
{Charge électrique\enspace·\enspace Loi de Coulomb\enspace·\enspace Gravitation\enspace·\enspace Champ électrostatique\enspace·\enspace Champ de gravitation}
\objectif{Interpréter des expériences d'électrisation et d'influence électrostatique.}
\objectif{Utiliser la loi de Coulomb et la loi de la gravitation universelle, y compris sous forme vectorielle.}
\objectif{Citer les analogies entre l'interaction électrostatique et l'interaction gravitationnelle.}
\objectif{Utiliser les expressions vectorielles des champs électrostatique et de gravitation.}
\objectif{Caractériser localement une ligne de champ ; exploiter une carte de champ.}
\prerequis{Charge électrique élémentaire, structure de l'atome (classe de seconde).}
\prerequis{Vecteurs, forces, poids $P=m\,g$ (classe de seconde) ; puissances de dix.}
\end{ouverture}
\essentiel
\essSection{1}{Charges électriques et interaction électrostatique}
\begin{definition}[Charge électrique]
La charge électrique d'un corps, en coulombs (C), est un multiple entier de la charge élémentaire $e=\qty{1,60e-19}{C}$ :
$q=\pm n\,e$. Un corps s'électrise par \textbf{frottement} ou par \textbf{contact}, en gagnant des électrons (il devient négatif) ou en
en perdant (il devient positif). Deux corps chargés de même signe se repoussent ; de signes opposés, ils s'attirent.
\end{definition}
\begin{propriete}[Influence électrostatique]
Un corps chargé approché d'un corps neutre déplace les charges de celui-ci : la face proche porte une charge de signe
opposé, la face éloignée une charge de même signe. L'attraction l'emporte, car les charges opposées sont plus proches :
un corps chargé attire les corps neutres légers.
\end{propriete}
\schema{11-influence}
\essSection{2}{La loi de Coulomb}
\begin{loi}[Loi de Coulomb]
Deux charges ponctuelles $q_\mathrm{A}$ et $q_\mathrm{B}$, séparées de la distance $d$, exercent l'une sur l'autre des forces de même valeur :
\[F=k\,\frac{|q_\mathrm{A}\,q_\mathrm{B}|}{d^2}\qquad\text{et}\qquad\vec{F}_{\mathrm{A}/\mathrm{B}}=k\,\frac{q_\mathrm{A}\,q_\mathrm{B}}{d^2}\,\vec{u}_{\mathrm{AB}},\qquad k=\qty{9,0e9}{N.m^2.C^{-2}},\]
où $\vec{u}_{\mathrm{AB}}$ est le vecteur unitaire dirigé de A vers B. Par exemple, deux charges de \qty{10}{nC} distantes de \qty{10}{cm} se
repoussent avec une force de \qty{9,0e-5}{N}.
\end{loi}
\schema{11-coulomb}
\essSection{3}{L'interaction gravitationnelle}
\begin{loi}[Loi de la gravitation universelle]
Deux corps à symétrie sphérique de masses $m_\mathrm{A}$ et $m_\mathrm{B}$, de centres distants de $d$, s'attirent :
\[F=G\,\frac{m_\mathrm{A}\,m_\mathrm{B}}{d^2}\qquad\text{et}\qquad\vec{F}_{\mathrm{A}/\mathrm{B}}=-G\,\frac{m_\mathrm{A}\,m_\mathrm{B}}{d^2}\,\vec{u}_{\mathrm{AB}},\qquad G=\qty{6,67e-11}{N.m^2.kg^{-2}}.\]
La Terre et la Lune s'attirent ainsi avec une force de \qty{1,98e20}{N}.
\end{loi}
\begin{propriete}[Analogies et différences]
\par\smallskip\noindent\hfill\begin{tabular}{lll}
\toprule
 & Interaction électrostatique & Interaction gravitationnelle \\
\midrule
Grandeur en jeu & charge $q$ (C) & masse $m$ (kg) \\
Dépendance avec la distance & en $\dfrac{1}{d^2}$ & en $\dfrac{1}{d^2}$ \\
Sens & attractive ou répulsive & toujours attractive \\
Domaine dominant & atomes, molécules, matière & astres, chute des corps \\
\bottomrule
\end{tabular}\hfill\null\par\smallskip
\end{propriete}
\essSection{4}{La notion de champ}
\begin{definition}[Champs électrostatique et de gravitation]
Un objet chargé crée en tout point de l'espace un \textbf{champ électrostatique} $\vec{E}$ (en \unit{V.m^{-1}} ou \unit{N.C^{-1}}) ; une charge $q$ placée
en ce point subit la force $\vec{F}=q\,\vec{E}$. De même, une masse crée un \textbf{champ de gravitation} $\vec{g}$ (en \unit{N.kg^{-1}}) ; une
masse $m$ y subit la force $\vec{F}=m\,\vec{g}$. Pour une charge ponctuelle $q$ et un astre de masse $M$, à la distance $d$ :
\[\vec{E}=k\,\frac{q}{d^2}\,\vec{u}\qquad\qquad\vec{g}=-G\,\frac{M}{d^2}\,\vec{u},\qquad\text{$\vec{u}$ étant dirigé de la source vers le point.}\]
À la surface de la Terre, $g=G\,\dfrac{M_\mathrm{T}}{R_\mathrm{T}^2}=\qty{9,81}{N.kg^{-1}}$.
\end{definition}
\begin{propriete}[Lignes de champ, champ uniforme]
Une \textbf{ligne de champ} est tangente en chacun de ses points au vecteur champ et orientée dans son sens. Les lignes sont
d'autant plus serrées que le champ est intense. Entre deux plaques parallèles soumises à la tension $U$ et distantes de $d$,
le champ est \textbf{uniforme} : il garde la même direction, le même sens et la même valeur $E=\dfrac{U}{d}$, de la plaque positive vers la
plaque négative.
\end{propriete}
\schema{11-champs}
\schema{11-champ-g}
\begin{python}
G, M, R = 6.67e-11, 5.97e24, 6.37e6        # SI : Terre
for h in [0, 400e3, 2000e3, 35786e3]:      # altitudes en m
    g = G * M / (R + h)**2
    print(round(h / 1000), 'km :', round(g, 2), 'N/kg')
\end{python}
\begin{console}
0 km : 9.81 N/kg
400 km : 8.69 N/kg
2000 km : 5.68 N/kg
35786 km : 0.22 N/kg
\end{console}
\noindent\capacite{Capacité numérique} Le programme calcule la valeur du champ de gravitation terrestre à différentes altitudes : la
distance au centre de la Terre vaut $R+h$. À l'altitude de la Station spatiale internationale, le champ vaut encore
\qty{89}{\percent} de sa valeur au sol.
\begin{propriete}[Quelques ordres de grandeur]
\par\smallskip\noindent\hfill\begin{tabular}{lr@{\qquad}lr}
\toprule
Champ de gravitation & \unit{N.kg^{-1}} & Champ électrostatique & \unit{V.m^{-1}} \\
\midrule
Surface de la Terre & \num{9,81} & Air au sol, par beau temps & $\approx10^2$ \\
Station spatiale (\qty{400}{km}) & \num{8,69} & Sol, sous un nuage d'orage & $\approx10^4$ \\
Surface de la Lune & \num{1,62} & Étincelle dans l'air & $\approx\num{3e6}$ \\
Orbite géostationnaire & \num{0,22} & Électron de l'atome d'hydrogène & $\approx\num{5e11}$ \\
\bottomrule
\end{tabular}\hfill\null\par\smallskip
\end{propriete}
\begin{attention}
Ne pas confondre une force et un champ : le champ existe en un point même en l'absence de charge ou de masse pour le
« ressentir ». Dans les lois de Coulomb et de Newton, $d$ est la distance entre les \emph{centres}, pas l'altitude ; et les
charges et distances s'expriment en unités du Système international (C et m).
\end{attention}
\begin{experience}[Visualiser des lignes de champ]
Dans une cuve contenant de l'huile de ricin, on saupoudre des grains de semoule et l'on plonge deux électrodes reliées à
un générateur de haute tension. Les grains s'alignent le long des lignes de champ : rayonnantes autour d'une électrode
ponctuelle, parallèles entre deux plaques.
\end{experience}
\begin{rediger}
Pour une force : préciser sa direction, son sens et sa valeur, et faire un schéma. Pour un calcul, écrire la loi
littérale, convertir les données en unités SI (nC en C, cm en m, km en m), puis calculer en notation scientifique.
\end{rediger}
\exercices

\rubrique{Découverte}
\begin{blocLivre}{ex-11-1}
\begin{exercice}[1]{Électriser une règle}
On frotte une règle en plastique avec un chiffon de laine : la règle acquiert une charge $q=\qty{-3,2}{nC}$.
\begin{donnees}
$e=\qty{1,60e-19}{C}$.
\end{donnees}
\begin{questions}
\item La règle a-t-elle gagné ou perdu des électrons ? Combien ?
\item Quelle est la charge du chiffon après le frottement ? Justifier.
\item La règle attire de petits morceaux de papier neutres. Expliquer.
\item Pourquoi une règle métallique tenue à la main ne s'électrise-t-elle pas de la même manière ?
\end{questions}
\end{exercice}
\end{blocLivre}
\begin{blocLivre}{ex-11-2}
\begin{exercice}[1]{Appliquer la loi de Coulomb}
\schema{11-ex02}
\begin{donnees}
$k=\qty{9,0e9}{N.m^2.C^{-2}}$.
\end{donnees}
\begin{questions}
\item Les charges s'attirent-elles ou se repoussent-elles ?
\item Calculer la valeur de la force exercée par A sur B.
\item Représenter les deux forces avec l'échelle \qty{1}{cm} pour \qty{1,0e-5}{N}.
\item Que devient la valeur de la force si la distance double ?
\end{questions}
\end{exercice}
\end{blocLivre}
\begin{blocLivre}{ex-11-3}
\begin{exercice}[1]{La Terre et la Lune}
\begin{donnees}
$G=\qty{6,67e-11}{N.m^2.kg^{-2}}$ ; $M_\mathrm{T}=\qty{5,97e24}{kg}$ ; $M_\mathrm{L}=\qty{7,35e22}{kg}$ ; $M_\mathrm{S}=\qty{1,99e30}{kg}$ ; distance Terre-Lune \qty{3,84e8}{m} ;
distance Terre-Soleil \qty{1,50e11}{m}.
\end{donnees}
\begin{questions}
\item Calculer la valeur de la force exercée par la Terre sur la Lune.
\item Quelle est la valeur de la force exercée par la Lune sur la Terre ? Représenter les deux forces sur un schéma.
\item Calculer la force exercée par le Soleil sur la Terre.
\item Combien de fois l'attraction du Soleil sur la Terre dépasse-t-elle celle de la Lune ?
\end{questions}
\end{exercice}
\end{blocLivre}
\begin{blocLivre}{ex-11-4}
\begin{exercice}[1]{Le champ de pesanteur terrestre}
\begin{donnees}
$G=\qty{6,67e-11}{N.m^2.kg^{-2}}$ ; $M_\mathrm{T}=\qty{5,97e24}{kg}$ ; $R_\mathrm{T}=\qty{6,37e6}{m}$. La Station spatiale internationale orbite à \qty{400}{km} d'altitude.
\end{donnees}
\begin{questions}
\item Calculer la valeur du champ de gravitation à la surface de la Terre.
\item Calculer le poids d'un astronaute de \qty{70}{kg} au sol.
\item Calculer le champ de gravitation à l'altitude de la station, puis la force exercée par la Terre sur l'astronaute.
\item L'astronaute « flotte » dans la station. Est-ce parce que la Terre ne l'attire plus ?
\end{questions}
\end{exercice}
\end{blocLivre}
\begin{blocLivre}{ex-11-5}
\begin{exercice}[1]{Lire une carte de champ}
\schema{11-ex05}
\begin{questions}
\item Indiquer le signe de la charge placée en A et de celle placée en B. Justifier.
\item Représenter, sans souci d'échelle, la direction et le sens du champ aux points M, N et P.
\item Représenter la force subie par une charge négative placée en chacun de ces points.
\item En quel point, de M ou de N, le champ est-il le plus intense ? Justifier.
\end{questions}
\end{exercice}
\end{blocLivre}
\rubrique{Entraînement}
\begin{blocLivre}{ex-11-6}
\begin{exercice}[2]{Le champ de gravitation en Python}
On reprend le programme de l'essentiel.
\begin{questions}
\item \capacite{Capacité numérique}Pourquoi le programme utilise-t-il \code{R + h} et non \code{h} ?
\item Prévoir la valeur affichée pour une altitude de \qty{10000}{km}.
\item À quelle altitude le champ est-il divisé par 4 par rapport au sol ? Justifier sans calcul numérique.
\item Modifier le programme pour calculer le champ à la surface de la Lune ($M=\qty{7,35e22}{kg}$, $R=\qty{1,74e6}{m}$).
\end{questions}
\end{exercice}
\end{blocLivre}
\begin{blocLivre}{ex-11-7}
\begin{exercice}[2]{L'atome d'hydrogène}
Dans le modèle de Bohr, l'électron de l'atome d'hydrogène tourne autour du proton à la distance $r=\qty{5,3e-11}{m}$.
\begin{donnees}
$e=\qty{1,60e-19}{C}$ ; $m_\mathrm{e}=\qty{9,11e-31}{kg}$ ; $m_\mathrm{p}=\qty{1,67e-27}{kg}$ ; $k=\qty{9,0e9}{SI}$ ; $G=\qty{6,67e-11}{SI}$.
\end{donnees}
\begin{questions}
\item Calculer la force électrostatique exercée par le proton sur l'électron. Est-elle attractive ou répulsive ?
\item Calculer la force de gravitation entre ces deux particules.
\item Calculer le rapport de ces deux forces. Quelle interaction assure la cohésion de l'atome ?
\item Pourquoi, à l'échelle des astres, est-ce au contraire la gravitation qui domine ?
\end{questions}
\end{exercice}
\end{blocLivre}
\begin{blocLivre}{ex-11-8}
\begin{exercice}[2]{Deux lois si semblables}
\begin{questions}
\item Écrire la loi de Coulomb et la loi de la gravitation universelle sous forme vectorielle.
\item Citer deux analogies entre ces lois.
\item Citer deux différences.
\item Expliquer pourquoi le signe « $-$ » apparaît dans l'expression vectorielle de la force de gravitation et pas dans
celle de la loi de Coulomb.
\end{questions}
\end{exercice}
\end{blocLivre}
\begin{blocLivre}{ex-11-9}
\begin{exercice}[2]{Un électron dans un condensateur}
\schema{11-ex09}
\begin{donnees}
$e=\qty{1,60e-19}{C}$ ; $m_\mathrm{e}=\qty{9,11e-31}{kg}$ ; $g=\qty{9,81}{N.kg^{-1}}$.
\end{donnees}
\begin{questions}
\item Calculer la valeur du champ électrostatique entre les plaques. Représenter quelques lignes de champ.
\item Calculer la valeur de la force électrostatique subie par l'électron et la représenter.
\item Calculer le poids de l'électron. Que peut-on conclure ?
\item Que se passerait-il pour un proton placé au même endroit ?
\end{questions}
\end{exercice}
\end{blocLivre}
\begin{blocLivre}{ex-11-10}
\begin{exercice}[2]{Le pendule électrostatique}
\schema{11-ex10}
\begin{questions}
\item Étape 1 : la boule est neutre. Expliquer pourquoi elle est attirée par la règle. Représenter la répartition des charges.
\item Étape 2 : que se passe-t-il au moment du contact ?
\item Étape 3 : expliquer pourquoi la boule est maintenant repoussée.
\item Que se passerait-il avec une règle chargée positivement ?
\end{questions}
\end{exercice}
\end{blocLivre}
\begin{blocLivre}{ex-11-11}
\begin{exercice}[2]{Entre la Terre et la Lune}
Un vaisseau va de la Terre à la Lune. Il existe sur le trajet un point P où les champs de gravitation de la Terre et de la
Lune se compensent.
\begin{donnees}
$M_\mathrm{T}=\qty{5,97e24}{kg}$ ; $M_\mathrm{L}=\qty{7,35e22}{kg}$ ; $d_\mathrm{TL}=\qty{3,84e8}{m}$.
\end{donnees}
\begin{questions}
\item Que dire des deux champs en P ? Faire un schéma.
\item On note $x$ la distance entre le centre de la Terre et P. Montrer que $\dfrac{M_\mathrm{T}}{x^2}=\dfrac{M_\mathrm{L}}{(d_\mathrm{TL}-x)^2}$.
\item En déduire que $\dfrac{x}{d_\mathrm{TL}-x}=\sqrt{\dfrac{M_\mathrm{T}}{M_\mathrm{L}}}$, puis calculer $x$.
\item Pourquoi ce point est-il beaucoup plus proche de la Lune que de la Terre ?
\end{questions}
\end{exercice}
\end{blocLivre}
\rubrique{Objectif DS}
\begin{blocLivre}{ex-11-12}
\begin{exercice}[3]{La photocopieuse}
\begin{docu}{1}{Le principe}
Dans une photocopieuse laser, un tambour est chargé électriquement aux endroits correspondant au texte. Des grains de
poudre, le toner, chargés de signe opposé, sont attirés par ces zones, puis transférés et fixés sur le papier.
\end{docu}
\begin{docu}{2}{Un grain de toner}
Un grain sphérique de diamètre \qty{10}{\micro m}, de masse volumique \qty{1,2e3}{kg.m^{-3}}, porte la charge $q=\qty{2,0e-14}{C}$. Près du tambour, le
champ électrostatique vaut environ \qty{1,0e6}{V.m^{-1}}.
\end{docu}
\begin{donnees}
$g=\qty{9,81}{N.kg^{-1}}$ ; $e=\qty{1,60e-19}{C}$ ; volume d'une sphère : $V=\dfrac{4}{3}\,\pi\,r^3$.
\end{donnees}
\begin{questions}
\item Combien de charges élémentaires porte un grain ?
\item Calculer la masse puis le poids d'un grain.
\item Calculer la force électrostatique subie par un grain près du tambour.
\item \initiative Comparer ces deux forces et conclure sur le fonctionnement de la photocopieuse.
\end{questions}
\end{exercice}
\end{blocLivre}
\begin{blocLivre}{ex-11-13}
\begin{exercice}[3]{Un satellite géostationnaire}
\begin{docu}{1}{Une orbite particulière}
Un satellite géostationnaire tourne autour de la Terre en 23~h~56~min, dans le plan de l'équateur, à \qty{35786}{km} d'altitude :
il paraît immobile vu du sol, ce qui permet de pointer les antennes une fois pour toutes.
\end{docu}
\begin{docu}{2}{Le satellite}
Masse : \qty{3000}{kg}.
\end{docu}
\begin{donnees}
$G=\qty{6,67e-11}{N.m^2.kg^{-2}}$ ; $M_\mathrm{T}=\qty{5,97e24}{kg}$ ; $R_\mathrm{T}=\qty{6,37e6}{m}$ ; $g_0=\qty{9,81}{N.kg^{-1}}$ au sol.
\end{donnees}
\begin{questions}
\item Calculer la distance entre le satellite et le centre de la Terre.
\item Calculer le champ de gravitation terrestre à cette distance.
\item Calculer la force exercée par la Terre sur le satellite et la comparer à son poids au sol.
\item \initiative Exprimer le champ à l'altitude $h$ en fonction de $g_0$, $R_\mathrm{T}$ et $h$, et vérifier le résultat de la question 2.
\end{questions}
\end{exercice}
\end{blocLivre}
\begin{blocLivre}{ex-11-14}
\begin{exercice}[3]{L'expérience de Millikan}
\begin{docu}{1}{Mesurer la charge élémentaire}
En 1909, Robert Millikan pulvérise de fines gouttelettes d'huile entre deux plaques horizontales. Certaines se chargent
par frottement. En réglant la tension entre les plaques, il immobilise une gouttelette : la force électrostatique compense
alors son poids. Il montre ainsi que toutes les charges sont des multiples d'une même charge.
\end{docu}
\begin{docu}{2}{Le dispositif}
\schema{11-ex14}
La gouttelette étudiée, chargée négativement, a pour masse \qty{2,94e-15}{kg}.
\end{docu}
\begin{donnees}
$g=\qty{9,81}{N.kg^{-1}}$ ; $e=\qty{1,60e-19}{C}$.
\end{donnees}
\begin{questions}
\item Faire l'inventaire des forces qui s'exercent sur la gouttelette immobile et les représenter.
\item Quelle plaque doit être chargée positivement ? Justifier.
\item Calculer le champ entre les plaques, puis la charge de la gouttelette.
\item \initiative Vérifier que cette charge est un multiple de la charge élémentaire.
\end{questions}
\end{exercice}
\end{blocLivre}
\begin{blocLivre}{ex-11-15}
\begin{exercice}[3]{Peser la Terre}
\begin{docu}{1}{Une idée de Cavendish}
En 1798, Henry Cavendish mesure la constante $G$ à l'aide d'une balance de torsion. On dit souvent qu'il a ainsi « pesé la
Terre » : connaissant $G$, le rayon terrestre et la valeur du champ de pesanteur, on peut en effet en déduire la masse de la
Terre.
\end{docu}
\begin{docu}{2}{Des roches}
Les roches de la croûte terrestre ont une masse volumique d'environ \qty{2,7e3}{kg.m^{-3}}.
\end{docu}
\begin{donnees}
$G=\qty{6,67e-11}{N.m^2.kg^{-2}}$ ; $R_\mathrm{T}=\qty{6,37e6}{m}$ ; $g=\qty{9,81}{N.kg^{-1}}$ ; volume d'une sphère : $\dfrac{4}{3}\,\pi\,R^3$.
\end{donnees}
\begin{questions}
\item Exprimer la masse de la Terre en fonction de $g$, $R_\mathrm{T}$ et $G$.
\item Calculer la masse de la Terre.
\item \initiative Calculer la masse volumique moyenne de la Terre et la comparer à celle des roches de surface.
\item Que peut-on en déduire sur l'intérieur de la Terre ?
\end{questions}
\end{exercice}
\end{blocLivre}
